\documentclass[../../../../main.tex]{subfiles}
\begin{document}

下記の表のように，次の規則によって数列を順に並べる．

\begin{figure}[htbp]
  \centering
  \begin{tikzpicture}[xscale=0.7, yscale=0.8]
    % 第1行
    \node at (-5, 3) {(第1行)};
    \node (A11) at (-1.5, 3) {$1$};
    \node (A12) at (-0.5, 3) {$1$};
    \node (A13) at (0.5, 3) {$1$};
    \node (A14) at (1.5, 3) {$1$};

    % 第2行
    \node at (-5, 2) {(第2行)};
    \node (A20) at (-2.5, 2) {$1$};
    \node (A21) at (-1, 2) {$2$};
    \node (A22) at (0, 2) {$2$};
    \node (A23) at (1, 2) {$2$};
    \node (A24) at (2.5, 2) {$1$};

    \draw (A11) -- (A21);
    \draw (A12) -- (A21);
    \draw (A12) -- (A22);
    \draw (A13) -- (A22);
    \draw (A13) -- (A23);
    \draw (A14) -- (A23);

    % 第3行
    \node at (-5, 1) {(第3行)};
    \node (A30) at (-3.5, 1) {$1$};
    \node (A31) at (-2, 1) {$3$};
    \node (A32) at (-0.8, 1) {$4$};
    \node (A33) at (0.8, 1) {$4$};
    \node (A34) at (2, 1) {$3$};
    \node (A35) at (3.5, 1) {$1$};

    \draw (A20) -- (A31);
    \draw (A21) -- (A31);
    \draw (A21) -- (A32);
    \draw (A22) -- (A32);
    \draw (A22) -- (A33);
    \draw (A23) -- (A33);
    \draw (A23) -- (A34);
    \draw (A24) -- (A34);

    % 第4行
    \node at (-5, 0) {(第4行)};
    \node (A40) at (-4.2, 0) {$1$};
    \node (A41) at (-2.8, 0) {$4$};
    \node (A42) at (-1.5, 0) {$7$};
    \node (A43) at (0, 0) {$8$};
    \node (A44) at (1.5, 0) {$7$};
    \node (A45) at (2.8, 0) {$4$};
    \node (A46) at (4.2, 0) {$1$};

    \draw (A30) -- (A41);
    \draw (A31) -- (A41);
    \draw (A31) -- (A42);
    \draw (A32) -- (A42);
    \draw (A32) -- (A43);
    \draw (A33) -- (A43);
    \draw (A33) -- (A44);
    \draw (A34) -- (A44);
    \draw (A34) -- (A45);
    \draw (A35) -- (A45);

    % 点々
    \node at (0, -0.8) {$\dots\dots\dots\dots\dots\dots\dots\dots$};
  \end{tikzpicture}
\end{figure}

\begin{enumerate}
  \item 第1行は1，1，1，1である．
  \item 第2行以下では，左右両端の数は1であり，その他の数は左上の数と右上の数との和である．
\end{enumerate}

第 $n$ 行の数列を ${}_nA_0,{}_nA_1,\cdots,{}_nA_k,\cdots,{}_nA_{n+2}$ とかけば
\begin{align*}
  (1+x^2){(1+x)}^n=\sum_{k=0}^{n+2}{}_nA_kx^k
\end{align*}
が成り立つことを示せ．

\end{document}
